\documentclass[10pt,a4paper]{amsart}
\usepackage[margin=19mm]{geometry}
\usepackage{amsmath,amssymb,mathtools}
\usepackage[scr=rsfs,cal=euler]{mathalfa}
\usepackage{xfrac}
\usepackage[unicode,hidelinks,bookmarks=false]{hyperref}
\newcommand{\ie}{i.e.\ }
\newcommand{\cf}{cf.\ }
\newcommand{\resp}{respectively\ }
\newcommand{\Subsection}{\subsection}
\providecommand{\nobreakdash}{\nobreak\hbox{-}}
\setlength{\emergencystretch}{2em}
\multlinegap0pt
\usepackage{mathtools}
\usepackage{xcolor}
\usepackage{amssymb}
\newtheorem{lemma}{Lemma}
\title{A generalization of Boppana's entropy inequality}
\author{Boon Suan Ho}
\date{Source: arXiv:2601.19327v1 (CC BY 4.0)}
\begin{document}
\maketitle

\noindent\emph{Source and licence.} A generalization of Boppana's entropy inequality. Version arXiv:2601.19327v1 (CC BY 4.0), distributed by arXiv under the Creative Commons Attribution 4.0 International licence, http://creativecommons.org/licenses/by/4.0/, as recorded in the arXiv licence metadata for this exact version and shown on the abstract page https://arxiv.org/abs/2601.19327v1. Full licence notice: \texttt{licenses/arxiv-2601.19327-CC-BY-4.0.txt}.\par\smallskip

\noindent\textit{Editorial scope.} Counted unit: the article's lemma lemma:l4 (source0.tex lines 199--201). The article's complete original proof of it is delivered with it (source0.tex lines 203--215, one \texttt{proof} environment of 13 lines). No mathematics is written, repaired or re-derived: the statement and the proof are the article's own lines, in the article's own notation, and no retained line is substituted or re-flowed.  No further source-local context is needed: every cross-reference of every retained line resolves inside the two delivered spans.  Nothing was omitted from any retained span, so the package carries no omission record.  No retained line contains a citation, so the wrapper carries no bibliography and no bibliography entry.  No retained line contains a figure environment, an image-inclusion command or drawing code, so no image is delivered and the package declares no asset dependency.  Editorial formatting only, in addition: an AMS \texttt{amsart} preamble that restates the article's own declarations and macros used by the retained lines, namely the environments \texttt{lemma} on separate counters, together with the standard packages those lines need.  The retained environments print in this document as Lemma 1 in order of appearance, whereas the article prints the same items with the numbering of its own full document.  The complete native source of the article as distributed in the arXiv e-print for this exact version is bundled unchanged as \texttt{sources/193450/source0.tex}.
\begin{lemma} \label{lemma:l4}
Let $U(x)=\log(x)\log(1-x)/h(x)$. Then $U$ is decreasing on $(0,1/2]$.
\end{lemma}

\begin{proof}
It suffices to prove that $1/U$ is increasing on $(0,1/2)$.
Since \begin{equation}
{1\over U(x)}={-x\over\log(1-x)}+{-(1-x)\over\log x}=L(1,1-x)+L(1,x)
\end{equation}
where $L(a,b)\coloneqq(a-b)/(\log a-\log b)$ is the \emph{logarithmic mean},
the identity $L(a,b)=\int_0^1a^{1-s}b^s\,ds$ yields
$f(t)\coloneqq L(1,t)=\int_0^1t^s\,ds$. 
Then $f''(t)=\int_0^1s(s-1)t^{s-2}\,ds<0$ on $(0,1)$, 
so $f(t)$ is strictly concave.
Thus $f'(t)$ is strictly decreasing, and we conclude that 
$(1/U(x))'=f'(x)-f'(1-x)>0$ for $x\in(0,1/2)$ as needed.
\end{proof}

\end{document}
